% fig-scramble.tex -- the scramble ablation and why its difference is leak-immune.
% No data dependency.
% Layout note: the prompt boxes are ~2.2cm tall once the shaded spans are set, which is far more
% than a four-line block looks like it should need. Every band below them is placed against that
% measured height, not against an estimate.
\begin{figure}[htbp]
\centering
\begin{tikzpicture}[
  x=1cm, y=1cm, font=\small,
  lbl/.style={font=\scriptsize, text=black!70},
  strong/.style={font=\scriptsize, text=black!90},
  promptbox/.style={draw=black!35, rounded corners=2pt, inner sep=5pt, fill=black!2},
  ar/.style={-{Stealth[length=4.5pt]}, black!65},
]
\newcommand{\spanhl}[1]{\colorbox{black!14}{#1}}

\node[lbl, anchor=west] at (0.15,0.00) {real prompt};
\node[promptbox, anchor=north west] (p1) at (0.15,-0.26) {%
  \renewcommand{\arraystretch}{0.9}%
  \begin{tabular}{@{}l@{~~}l@{}}
  \ttfamily\scriptsize Cell line:    & \ttfamily\scriptsize MCF7 \\
  \ttfamily\scriptsize Drug:         & \ttfamily\scriptsize \spanhl{Lapatinib} \\
  \ttfamily\scriptsize Mechanism:    & \ttfamily\scriptsize \spanhl{EGFR inhibitor} \\
  \ttfamily\scriptsize Control cell: & \ttfamily\scriptsize ACTB GAPDH \ldots \\
  \end{tabular}};

\node[lbl, anchor=west] at (0.15,-2.86) {scrambled prompt --- only the two shaded spans differ};
\node[promptbox, anchor=north west] (p2) at (0.15,-3.12) {%
  \renewcommand{\arraystretch}{0.9}%
  \begin{tabular}{@{}l@{~~}l@{}}
  \ttfamily\scriptsize Cell line:    & \ttfamily\scriptsize MCF7 \\
  \ttfamily\scriptsize Drug:         & \ttfamily\scriptsize \spanhl{Vorinostat} \\
  \ttfamily\scriptsize Mechanism:    & \ttfamily\scriptsize \spanhl{HDAC inhibitor} \\
  \ttfamily\scriptsize Control cell: & \ttfamily\scriptsize ACTB GAPDH \ldots \\
  \end{tabular}};

\node[draw=black!35, rounded corners=2pt, inner sep=6pt, align=center, fill=black!3]
  (truth) at (9.30,-2.55) {truth for\\\textbf{Lapatinib}\\{\scriptsize drawn once}};

\draw[ar] (p1.east) .. controls (6.6,-1.35) and (7.6,-1.85) .. (truth.north west);
\draw[ar] (p2.east) .. controls (6.6,-4.20) and (7.6,-3.30) .. (truth.south west);

\node[anchor=north, font=\small] at (9.30,-3.70)
  {$\gap = \NIR_{\text{real}} - \NIR_{\text{scrambled}}$};

\draw[ar, dashed, black!45] (1.20,-6.10) -- (8.10,-6.10);
\node[strong, anchor=south, align=center] at (4.65,-6.02)
  {the shared control cell enters \emph{both} arms, so control- and batch-derived};
\node[strong, anchor=north, align=center] at (4.65,-6.20)
  {leakage cancels in the difference};

\end{tikzpicture}
\caption[The scramble ablation]{\textbf{The scramble replaces the drug name and its mechanism and
nothing else.} Both arms are conditioned on the same control cell and scored against the same truth,
so any leakage through the control cell or the plate enters both arms in equal measure and cancels in
the difference $\gap$. This is why $\gap$, rather than a comparison against a drug-agnostic
predictor, is the quantity every drug-use claim in this thesis is anchored to: the drug-agnostic
comparison is not leak-immune, and a drug-agnostic mean baseline scores $\NIR = 0.399$ when the
comparison set may span plates against $0.161$ when it may not, precisely because it inherits that
leakage. A refinement matters here and is
developed in Section~\ref{sec:methods-controls}: if the swapped-in drug happens to have a
\emph{similar} true response then an unchanged output is correct rather than blind, so the swap
partner is chosen from three defined similarity strata and the evidence is the \emph{gradient} across
them rather than any single value.}
\label{fig:scramble}
\end{figure}
